\documentclass[aspectratio=169]{beamer}
\usepackage{tikz}
\usepackage{graphicx}
\usepackage{listings}
\usepackage{xcolor}
\usepackage{hyperref}
\usepackage[spanish,es-noquoting]{babel}
\usepackage[round]{natbib}

\lstset{
  basicstyle=\ttfamily\small,
  keywordstyle=\color{blue},
  stringstyle=\color{red},
  commentstyle=\color{gray},
  showstringspaces=false
}
% Minimal look
\usetheme{Warsaw}
%\setbeamertemplate{headline}
\setbeamertemplate{footline}{
  \leavevmode%
  \hbox{%
    % Left: Section name
    \begin{beamercolorbox}[wd=.33\paperwidth,ht=2.25ex,dp=1ex,left]{author in head/foot}%
      \hspace*{2ex}\insertsection
    \end{beamercolorbox}%
    % Center: Presentation title
    \begin{beamercolorbox}[wd=.34\paperwidth,ht=2.25ex,dp=1ex,center]{title in head/foot}%
      \insertshorttitle
    \end{beamercolorbox}%
    % Right: Slide numbers
    \begin{beamercolorbox}[wd=.33\paperwidth,ht=2.25ex,dp=1ex,right]{date in head/foot}%
      \insertframenumber{} / \inserttotalframenumber\hspace*{2ex}
    \end{beamercolorbox}%
  }%
  \vskip0pt%
}

\usecolortheme{default}
\usefonttheme{serif}
\graphicspath{{../}{./}}
\setbeamertemplate{caption}[numbered]

% Title info
\title[Module 5: Team Management and ]{Module 5: Team Management and Risk}
\subtitle{Red Teaming and Feedback}
\author{Ricardo Reyes-Jimenez}
\institute{Capgemini}
\date{\today}



\begin{document}

% Title slide
\begin{frame}
  \titlepage
\end{frame}

% Reset background for normal slides
\usebackgroundtemplate{}



\begin{frame}[fragile]{Session Objective}


\textbf{By the end of this module, the Engagement Manager will be able to:}

1. Apply the \textbf{Red Teaming} technique to stress-test project plans and timelines.
2. Anticipate operational bottlenecks before they occur.
3. Draft objective, constructive performance feedback for delivery team members.

\end{frame}



\begin{frame}[fragile]{1. What is Red Teaming?}


\textbf{Playing the Devil's Advocate to save your project.}

\begin{itemize}
  \item \textbf{The Concept:} In cybersecurity, the ``Red Team'' attacks the system to find vulnerabilities. For an EM, Red Teaming means asking the AI to find the flaws in your project plan.
  \item \textbf{Why it works:} Human bias makes us overly optimistic about our own timelines. AI can objectively point out unrealistic dependencies \cite{redteam2026}.
\end{itemize}
\end{frame}



\begin{frame}[fragile]{Exercise 1: Plan Stress-Test (Red Teaming)}


\textbf{Context:} You have drafted a 4-week aggressive timeline for a software rollout. You need to know where it's most likely to fail.
\textbf{Task:} Ask Copilot to attack your plan.

\begin{itemize}
  \item \textbf{Draft Prompt to Try:}
\end{itemize}
\textit{``Act as a highly critical, experienced Agile Delivery Director. Review my proposed 4-week rollout plan: [Paste Timeline]. Perform a Red Team analysis. Identify the 3 most likely points of failure, resource bottlenecks, or unrealistic assumptions. Be brutally honest and suggest a mitigation tactic for each flaw.''}

\end{frame}



\begin{frame}[fragile]{2. Managing the Delivery Team}


\textbf{Constructive feedback is hard; emotional feedback is destructive.}

\begin{itemize}
  \item As an EM, you must often deliver difficult feedback to team members who are underperforming or missing deadlines.
  \item MS Copilot can help you strip the emotion from your notes and reframe them into actionable, objective guidance.
\end{itemize}
\end{frame}



\begin{frame}[fragile]{Exercise 2: Constructive Feedback Generation}


\textbf{Context:} A senior developer has been missing daily standups and delivering code late, causing delays for the QA team. You are frustrated.
\textbf{Task:} Use Copilot to draft a professional feedback script.

\begin{itemize}
  \item \textbf{Draft Prompt to Try:}
\end{itemize}
*"I need to give feedback to a Senior Developer. My raw, frustrated notes are: 'He keeps skipping standups, his code is late, and QA is blocked because of him. He needs to wake up.'
Rewrite this as a script for a 1-on-1 meeting. Remove the emotion. Focus on the objective impact on the project timeline and the QA team. Provide a constructive opening and a clear call to action for improvement."*

\end{frame}



\begin{frame}[fragile]{Recap \& Next Session}


\textbf{Today we learned:}
\begin{itemize}
  \item How to use AI as a ``Red Team'' to objectively find flaws in our project plans.
  \item The value of AI in neutralizing emotion when dealing with team performance issues.
  \item Translating frustration into actionable, professional feedback.
\end{itemize}
\textbf{Next Session:}
\textit{Module 6: Advanced Prompting Techniques.} We will learn how to clone your personal writing style (Few-Shot) and how to force the AI to reason sequentially (Chain-of-Thought) for complex SteerCo reports.
\end{frame}




\begin{frame}[allowframebreaks]{References}
  \bibliographystyle{apalike}
  \bibliography{references}
\end{frame}


\end{document}